% Options for packages loaded elsewhere
% Options for packages loaded elsewhere
\PassOptionsToPackage{unicode}{hyperref}
\PassOptionsToPackage{hyphens}{url}
\PassOptionsToPackage{dvipsnames,svgnames,x11names}{xcolor}
%
\documentclass[
  letterpaper,
  DIV=11,
  numbers=noendperiod]{scrartcl}
\usepackage{xcolor}
\usepackage[margin=1in]{geometry}
\usepackage{amsmath,amssymb}
\setcounter{secnumdepth}{-\maxdimen} % remove section numbering
\usepackage{iftex}
\ifPDFTeX
  \usepackage[T1]{fontenc}
  \usepackage[utf8]{inputenc}
  \usepackage{textcomp} % provide euro and other symbols
\else % if luatex or xetex
  \usepackage{unicode-math} % this also loads fontspec
  \defaultfontfeatures{Scale=MatchLowercase}
  \defaultfontfeatures[\rmfamily]{Ligatures=TeX,Scale=1}
\fi
\usepackage{lmodern}
\ifPDFTeX\else
  % xetex/luatex font selection
\fi
% Use upquote if available, for straight quotes in verbatim environments
\IfFileExists{upquote.sty}{\usepackage{upquote}}{}
\IfFileExists{microtype.sty}{% use microtype if available
  \usepackage[]{microtype}
  \UseMicrotypeSet[protrusion]{basicmath} % disable protrusion for tt fonts
}{}
\makeatletter
\@ifundefined{KOMAClassName}{% if non-KOMA class
  \IfFileExists{parskip.sty}{%
    \usepackage{parskip}
  }{% else
    \setlength{\parindent}{0pt}
    \setlength{\parskip}{6pt plus 2pt minus 1pt}}
}{% if KOMA class
  \KOMAoptions{parskip=half}}
\makeatother
% Make \paragraph and \subparagraph free-standing
\makeatletter
\ifx\paragraph\undefined\else
  \let\oldparagraph\paragraph
  \renewcommand{\paragraph}{
    \@ifstar
      \xxxParagraphStar
      \xxxParagraphNoStar
  }
  \newcommand{\xxxParagraphStar}[1]{\oldparagraph*{#1}\mbox{}}
  \newcommand{\xxxParagraphNoStar}[1]{\oldparagraph{#1}\mbox{}}
\fi
\ifx\subparagraph\undefined\else
  \let\oldsubparagraph\subparagraph
  \renewcommand{\subparagraph}{
    \@ifstar
      \xxxSubParagraphStar
      \xxxSubParagraphNoStar
  }
  \newcommand{\xxxSubParagraphStar}[1]{\oldsubparagraph*{#1}\mbox{}}
  \newcommand{\xxxSubParagraphNoStar}[1]{\oldsubparagraph{#1}\mbox{}}
\fi
\makeatother


\usepackage{longtable,booktabs,array}
\newcounter{none} % for unnumbered tables
\usepackage{calc} % for calculating minipage widths
% Correct order of tables after \paragraph or \subparagraph
\usepackage{etoolbox}
\makeatletter
\patchcmd\longtable{\par}{\if@noskipsec\mbox{}\fi\par}{}{}
\makeatother
% Allow footnotes in longtable head/foot
\IfFileExists{footnotehyper.sty}{\usepackage{footnotehyper}}{\usepackage{footnote}}
\makesavenoteenv{longtable}
\usepackage{graphicx}
\makeatletter
\newsavebox\pandoc@box
\newcommand*\pandocbounded[1]{% scales image to fit in text height/width
  \sbox\pandoc@box{#1}%
  \Gscale@div\@tempa{\textheight}{\dimexpr\ht\pandoc@box+\dp\pandoc@box\relax}%
  \Gscale@div\@tempb{\linewidth}{\wd\pandoc@box}%
  \ifdim\@tempb\p@<\@tempa\p@\let\@tempa\@tempb\fi% select the smaller of both
  \ifdim\@tempa\p@<\p@\scalebox{\@tempa}{\usebox\pandoc@box}%
  \else\usebox{\pandoc@box}%
  \fi%
}
% Set default figure placement to htbp
\def\fps@figure{htbp}
\makeatother





\setlength{\emergencystretch}{3em} % prevent overfull lines

\providecommand{\tightlist}{%
  \setlength{\itemsep}{0pt}\setlength{\parskip}{0pt}}



 


\KOMAoption{captions}{tableheading}
\makeatletter
\@ifpackageloaded{tcolorbox}{}{\usepackage[skins,breakable]{tcolorbox}}
\@ifpackageloaded{fontawesome5}{}{\usepackage{fontawesome5}}
\definecolor{quarto-callout-color}{HTML}{909090}
\definecolor{quarto-callout-note-color}{HTML}{0758E5}
\definecolor{quarto-callout-important-color}{HTML}{CC1914}
\definecolor{quarto-callout-warning-color}{HTML}{EB9113}
\definecolor{quarto-callout-tip-color}{HTML}{00A047}
\definecolor{quarto-callout-caution-color}{HTML}{FC5300}
\definecolor{quarto-callout-color-frame}{HTML}{acacac}
\definecolor{quarto-callout-note-color-frame}{HTML}{4582ec}
\definecolor{quarto-callout-important-color-frame}{HTML}{d9534f}
\definecolor{quarto-callout-warning-color-frame}{HTML}{f0ad4e}
\definecolor{quarto-callout-tip-color-frame}{HTML}{02b875}
\definecolor{quarto-callout-caution-color-frame}{HTML}{fd7e14}
\makeatother
\makeatletter
\@ifpackageloaded{caption}{}{\usepackage{caption}}
\AtBeginDocument{%
\ifdefined\contentsname
  \renewcommand*\contentsname{Table of contents}
\else
  \newcommand\contentsname{Table of contents}
\fi
\ifdefined\listfigurename
  \renewcommand*\listfigurename{List of Figures}
\else
  \newcommand\listfigurename{List of Figures}
\fi
\ifdefined\listtablename
  \renewcommand*\listtablename{List of Tables}
\else
  \newcommand\listtablename{List of Tables}
\fi
\ifdefined\figurename
  \renewcommand*\figurename{Figure}
\else
  \newcommand\figurename{Figure}
\fi
\ifdefined\tablename
  \renewcommand*\tablename{Table}
\else
  \newcommand\tablename{Table}
\fi
}
\@ifpackageloaded{float}{}{\usepackage{float}}
\floatstyle{ruled}
\@ifundefined{c@chapter}{\newfloat{codelisting}{h}{lop}}{\newfloat{codelisting}{h}{lop}[chapter]}
\floatname{codelisting}{Listing}
\newcommand*\listoflistings{\listof{codelisting}{List of Listings}}
\makeatother
\makeatletter
\makeatother
\makeatletter
\@ifpackageloaded{caption}{}{\usepackage{caption}}
\@ifpackageloaded{subcaption}{}{\usepackage{subcaption}}
\makeatother
\usepackage{bookmark}
\IfFileExists{xurl.sty}{\usepackage{xurl}}{} % add URL line breaks if available
\urlstyle{same}
\hypersetup{
  pdftitle={EFB 390: Wildlife Ecology and Management},
  pdfauthor={Dr.~Elie Gurarie},
  colorlinks=true,
  linkcolor={blue},
  filecolor={Maroon},
  citecolor={Blue},
  urlcolor={Blue},
  pdfcreator={LaTeX via pandoc}}


\title{EFB 390: Wildlife Ecology and Management}
\usepackage{etoolbox}
\makeatletter
\providecommand{\subtitle}[1]{% add subtitle to \maketitle
  \apptocmd{\@title}{\par {\large #1 \par}}{}{}
}
\makeatother
\subtitle{Syllabus --- Fall 2026}
\author{Dr.~Elie Gurarie}
\date{}
\begin{document}
\maketitle


\begin{tcolorbox}[enhanced jigsaw, arc=.35mm, bottomrule=.15mm, breakable, colback=white, colframe=quarto-callout-note-color-frame, left=2mm, leftrule=.75mm, opacityback=0, rightrule=.15mm, toprule=.15mm]
\begin{minipage}[t]{5.5mm}
\textcolor{quarto-callout-note-color}{\faInfo}
\end{minipage}%
\begin{minipage}[t]{\textwidth - 5.5mm}

\textbf{Lecture:} Tues \& Thurs 2:00--3:20, Illick 5

\textbf{Recitations:} Tues 3:30--4:25 (Baker 310), Tues 5:00--5:55
(Baker 314), Wed 3:45--4:40 (Baker 314), Thurs 8:00--8:55 (Baker 314)

\end{minipage}%
\end{tcolorbox}

\begin{itemize}
\item
  \textbf{Instructor}: \href{https://eligurarie.github.io/}{Dr.~Elie
  Gurarie}
\item
  \textbf{Office}: Illick 206
\item
  \textbf{Office hours}: Thursday 3:30--4:30 and by appointment
\item
  \textbf{Email:}
  \href{mailto:egurarie@esf.edu}{\nolinkurl{egurarie@esf.edu}}
\item
  \textbf{TAs}: Seanna Jobe
  (\href{mailto:sjobe@esf.edu}{\nolinkurl{sjobe@esf.edu}}) and Dani
  Venlen (\href{mailto:dverlen@esf.edu}{\nolinkurl{dverlen@esf.edu}})
\end{itemize}

Office hours and locations TBA.

\subsection{Course Description}\label{course-description}

This is a broad, foundational course. The overarching goal is to make
students familiar with fundamental topics in wildlife ecology and
management.

Wildlife \emph{ecology} is an extremely complex science, that explores
themes like population dynamics, behaviors, space use, disease, habitat,
trophic interactions, that is studied with a suite of rapidly evolving
tools --- field observations, advancing technology, statistics and
modeling.

Wildlife \emph{management} places all the complexity of wildlife ecology
into a sloppy social, political, historical, ethical and legal realm.
Most professional ``wildlife'' managers will openly admit that they
spend much more time trying to manage \emph{people}.

In this course, we will introduce methods, theories, concepts, and
contemporary research topics in wildlife ecology, placing these into the
very human inflected context of management.

We will hopefully also get you even more interested in wildlife ecology
than you maybe already are! It is my opinion that there is no better or
more interesting job than being a wildlife ecologist.

\subsection{Course Learning Objectives /
Outcomes}\label{course-learning-objectives-outcomes}

\begin{itemize}
\tightlist
\item
  Understand fundamental concepts in wildlife ecology and ecological
  theory, such as population dynamics, wildlife-habitat relationships,
  limiting factors, surveying and estimation, wildlife statistics.
\item
  Develop and hone research skills --- especially building a
  bibliography, understanding and synthesizing peer reviewed literature
  and grey literature on wildlife ecology and management.
\item
  Develop some rudimentary skills in R programming and statistical tools
  for fitting and interpreting commonly used models in wildlife ecology.
\item
  Learn to establish and maintain a camera trap study.
\item
  Become familiar with the historical context and current (conflicting)
  philosophies of wildlife management, how these apply to past, present,
  and future approaches to wildlife research and management.
\item
  Become familiar with the human dimensions of wildlife ecology and
  management including the role of social, economic, and political
  dimensions through analysis and synthesis of research articles,
  popular press articles, reports, and other media, and discussions.
\item
  Synthesize concepts in wildlife ecology, wildlife and habitat
  management, policy, and human dimensions components to describe,
  understand, and consider solutions to contemporary wildlife
  conservation and management challenges, using evidence-based
  approaches.
\end{itemize}

\subsection{College Learning Outcomes}\label{college-learning-outcomes}

Scientific Reasoning \textbar{} Quantitative Reasoning \textbar{}
Communication Skills

Technological and Information Literacy (Zotero!) \textbar{} Basic
programming and statistics (R!)

Values, Ethics and Diverse Perspectives \textbar{} Critical Thinking
\textbar{} Collaborative work

\subsection{Textbooks and Supplies}\label{textbooks-and-supplies}

There is no required textbook in this class, though a few texts can be
recommended. There will, however, be many readings assigned --- book
chapters, scientific articles, population articles, ``grey literature,''
some seminal readings. All materials will be made available on the
course Blackboard page. All lectures (and other important materials)
will be on the course website:
\url{https://eligurarie.github.io/EFB390/}. Bookmark this page!

\subsection{Attendance Policy}\label{attendance-policy}

This course will \textbf{primarily} cover material that does not appear
in your assigned readings or other out-of-class materials. Thus, you
should make an effort to attend all classes. If you must be absent for a
class, consult with me beforehand and make sure to obtain class notes
from a classmate.

We strongly encourage you to make sure to come to class prepared. If you
come to class consistently, come prepared having completed any assigned
reading or viewing, and pay attention, you are highly likely to succeed
in learning the concepts and knowledge I am hoping for you to learn and
earning a grade you are satisfied with.

In view of the AI policy (see below), there will be occasional
assessments or quizzes in recitation, for which attendance will be
obligatory.

Also --- visit us, all of us, at office hours (or at other times)! We
love to talk about this stuff.

\subsection{Grading Structure}\label{grading-structure}

Your grade will be based on 1 (or 2) quizzes, weekly assignments
(including several group projects during the semester), and a final
major project + final group presentation.

{\def\LTcaptype{none} % do not increment counter
\begin{longtable}[]{@{}ll@{}}
\toprule\noalign{}
Assignment & Grade Percentage \\
\midrule\noalign{}
\endhead
\bottomrule\noalign{}
\endlastfoot
Weekly Assignments & 30\% \\
Exams (2) & 25\% \\
Final project + presentation & 25\% \\
Participation + engagement + quizzes & 20\% \\
\textbf{\emph{Total}} & \textbf{\emph{100\%}} \\
\end{longtable}
}

Academic dishonesty, non-discrimination and inclusive excellence, and
others, can be found in the ESF student handbook
(\url{https://www.esf.edu/students/handbook/}).

\subsection{Policy on the Use of AI
Tools}\label{policy-on-the-use-of-ai-tools}

Large Language Models (LLMs) and other generative AI programs (ChatGPT,
Claude, Gemini, etc.) are becoming increasingly powerful and accessible.
As a practicing wildlife ecologist I use them routinely in my work to
speed up certain tasks, though not without having to learn how to use
those tools. However, there are real risks to relying on these tools too
much. Remember, and think of, AI is a \emph{tool}, not a substitute for
thought. As an analogy, we are, mainly, no longer hunters and gatherers
and could theoretically live a lifetime without exercising any muscles
at all. And yet --- we chose to move. In a similar way, \emph{choose to
use your brain!} Yes --- you can get a result with a single prompt. It
is much, much harder to internalize information and process ideas when
you don't work through them and sit on them and sleep on them and
discuss them before --- finally, going through the act of articulating
and communicating those ideas through writing. With that in mind, the
policy in this course will be to:

\begin{enumerate}
\def\labelenumi{\arabic{enumi}.}
\tightlist
\item
  \textbf{Be transparent}: If you use a generative AI in completing
  assignments, you must clearly disclose (a) which you used, and (b)
  what for.
\item
  \textbf{Have higher expectations}: Because these tools can accelerate
  some tasks, the standards for depth, accuracy, and originality in your
  work will be higher than in past versions of this course.
\item
  \textbf{More in-class assessments}: To ensure each student is
  developing their own skills, a slightly greater portion of your grade
  will come from in-class assessments (e.g., short quizzes during
  recitation).
\end{enumerate}

This is a dynamic and rapidly evolving aspect of teaching and learning
--- we can work together to make the most of these new opportunities.

\subsection{\texorpdfstring{\emph{Tentative} Class
Schedule}{Tentative Class Schedule}}\label{tentative-class-schedule}

Classes begin Tuesday, August 25. See the
\hyperref[academic-calendar]{academic calendar} below for holidays and
exam dates.

\emph{GL = Guest Lecture, filled out as speakers are confirmed.
Recitation topics listed in italics.}

{\def\LTcaptype{none} % do not increment counter
\begin{longtable}[]{@{}
  >{\raggedright\arraybackslash}p{(\linewidth - 6\tabcolsep) * \real{0.2500}}
  >{\raggedright\arraybackslash}p{(\linewidth - 6\tabcolsep) * \real{0.2500}}
  >{\raggedright\arraybackslash}p{(\linewidth - 6\tabcolsep) * \real{0.2500}}
  >{\raggedright\arraybackslash}p{(\linewidth - 6\tabcolsep) * \real{0.2500}}@{}}
\toprule\noalign{}
\begin{minipage}[b]{\linewidth}\raggedright
Module
\end{minipage} & \begin{minipage}[b]{\linewidth}\raggedright
Date
\end{minipage} & \begin{minipage}[b]{\linewidth}\raggedright
Day
\end{minipage} & \begin{minipage}[b]{\linewidth}\raggedright
Topic
\end{minipage} \\
\midrule\noalign{}
\endhead
\bottomrule\noalign{}
\endlastfoot
\textbf{I. Background} & Aug 25 & Tues. & Basics and Definitions \\
& Aug 25 & \emph{Recitation} & \emph{Wildlife in the News} \\
& Aug 27 & Thurs. & Library Science and Research (GL: \emph{Moon
Library}) \\
& Sep 1 & Tues. & Human-Wildlife History -- Part I \\
& Sep 1 & \emph{Recitation} & \emph{Referencing and Bibliographies} \\
& Sep 3 & Thurs. & Human-Wildlife History -- Part II \\
\textbf{II. Estimation and Sampling} & Sep 8 & Tues. & Counting Animals
-- Part I \\
& Sep 8 & \emph{Recitation} & \emph{Survey Methods + Camera Trap
Deployment} \\
& Sep 10 & Thurs. & Counting Animals -- Part II \\
& Sep 15 & Tues. & Counting Animals -- Sampling Uncertainty \\
& Sep 15 & \emph{Recitation} & \emph{Strategizing Flag Abundance
Estimation} \\
& Sep 17 & Thurs. & Flag Abundance -- Field Simulation + Data Entry \\
\textbf{III. Ranges and Habitats} & Sep 22 & Tues. & Mark-Recapture +
Index Counts (GL: \textbf{Irina Trukhanova}) \\
& Sep 22 & \emph{Recitation} & \emph{Count data analysis (R lab I)} \\
& Sep 24 & Thurs. & \emph{(no EG)} Distributions and Ranges \textbar{}
Niches and Habitats \\
& Sep 29 & Tues. & Modeling just about anything I \\
& Sep 29 & \emph{Recitation} & \emph{Wildlife Statistics (R lab II)} \\
& Oct 1 & Thurs. & Modeling just about anything II \\
& Oct 6 & Tues. & Forest Management for Wildlife (GL: \emph{Andrew
Vander Yacht}) \\
& Oct 6 & \emph{Recitation} & \emph{AIC Symposium Prep} \\
& Oct 8 & Thurs. & Mega symposium of mini-talks on delta AIC in Wildlife
Ecology \\
& Oct 13 & Tues. & \textbf{NO CLASS \textbar{} NO RECITATION} (Fall
Break) \\
& Oct 15 & Thurs. & Exam 1 Review \\
& Oct 20 & Tues. & \textbf{Exam I} \\
& Oct 20 & \emph{Recitation} & \emph{ID'ing, Storing and Managing Camera
Trap Data} \\
\textbf{IV. Populations} & Oct 22 & Thurs. & \emph{(no EG)} Population
Growth \\
& Oct 27 & Tues. & Limits to Growth \\
& Oct 27 & \emph{Recitation} & \emph{Estimating Population Growth (R Lab
3)} \\
& Oct 29 & Thurs. & Population Structure II (\emph{GL: \textbf{TBD}}) \\
& Nov 3 & Tues. & Species Interactions \textbar{} Wildlife Diseases \\
& Nov 3 & \emph{Recitation} & \emph{Structured Populations (R Lab 4)} \\
\textbf{V. Models of Management}* & Nov 5 & Thurs. & North American
Model of Wildlife Conservation \\
& Nov 10 & Tues. & Critiques of NAM (\emph{GL: \textbf{TBD}}) \\
& Nov 12 & Thurs. & Hunter Perspectives (\emph{GL: \textbf{TBD}}) \\
\textbf{VI. Wildlife Management in Practice} & Nov 17 & Tues. & Wildfowl
Management (GL: \textbf{\emph{Michael Schummer}}) \\
& Nov 19 & Thurs. & Wildlife Management Elsewhere (GL:
\textbf{\emph{TBD}}) \\
& Nov 24--26 & Tues./Thurs. & \textbf{NO CLASS \textbar{} NO RECITATION}
(Thanksgiving Recess) \\
& Dec 1 & Tues. & \emph{Guest Lecture TBD} (Physiology \textbar{} Feds
\textbar{} Collections) \\
& Dec 3 & Thurs. & Exam Review \\
& Dec 8 & Tues. & \textbf{Exam 2} (last day of classes) \\
\textbf{Finals} & Dec 10--15 & (TBD) & \textbf{Final Project
Presentations} \\
\end{longtable}
}

* \emph{For the final portion of class, recitations will be devoted to
guided progress on final projects.}

\begin{tcolorbox}[enhanced jigsaw, arc=.35mm, bottomrule=.15mm, bottomtitle=1mm, breakable, colback=white, colbacktitle=quarto-callout-important-color!10!white, colframe=quarto-callout-important-color-frame, coltitle=black, left=2mm, leftrule=.75mm, opacityback=0, opacitybacktitle=0.6, rightrule=.15mm, title=\textcolor{quarto-callout-important-color}{\faExclamation}\hspace{0.5em}{Important}, titlerule=0mm, toprule=.15mm, toptitle=1mm]

I retain the right to change anything at any time! In practice this is
nearly always to the benefit of the students.

\end{tcolorbox}

\subsubsection{Academic Calendar}\label{academic-calendar}

{\def\LTcaptype{none} % do not increment counter
\begin{longtable}[]{@{}ll@{}}
\toprule\noalign{}
Event & Date(s) \\
\midrule\noalign{}
\endhead
\bottomrule\noalign{}
\endlastfoot
Classes begin & August 24 \\
Labor Day (no classes) & September 7 \\
Fall Break (no classes) & October 12--13 \\
Mid-term grading & October 12--16 \\
Thanksgiving Recess (no classes) & November 22--29 \\
Last day of classes & December 8 \\
Final Exams & December 10--11 \\
Final Exams & December 14--15 \\
Grades due & December 22 \\
\end{longtable}
}

\subsection{Additional Important
Information}\label{additional-important-information}

\textbf{Religious holy days:} Students who miss coursework due to the
observance of a religious holy day will be given the opportunity to
complete the work missed within a reasonable time after the absence,
provided that the instructor is notified in advance (notify the course
instructor at least 2 weeks prior to the class or an exam that will be
missed).

\textbf{Scholastic dishonesty:} Students must act with integrity in
accordance to ESF's Code of Academic Integrity.

\textbf{Common courtesy:} Turn cell phones off, put on silent mode, or
whatever it takes to keep them quiet. No texting, emailing, etc. during
lecture. Please be on time.

\textbf{Disability Services:} SUNY-ESF works with the Office of
Disability Services (ODS) at Syracuse University, who is responsible for
coordinating disability-related accommodations. Students can contact ODS
at 804 University Avenue-Room 309, 315-443-4498 to schedule an
appointment and discuss their needs and the process for requesting
accommodations. Students may also contact the ESF Office of Student
Affairs, 110 Bray Hall, 315-470-6660 for assistance with the process. To
learn more about ODS, visit \url{http://disabilityservices.syr.edu}.
Authorized accommodation forms must be in the instructor's possession
one week prior to any anticipated accommodation. Since accommodations
may require early planning and generally are not provided retroactively,
please contact ODS as soon as possible.

\textbf{Diversity and Inclusion:} SUNY-ESF values diversity and
inclusion; we are committed to a climate of mutual respect and full
participation. Our goal is to create learning environments that are
usable, equitable, inclusive and welcoming. If there are aspects of the
instruction or design of this course that result in barriers to your
inclusion or accurate assessment or achievement, we invite any student
to meet with us to discuss additional strategies beyond accommodations
that may be helpful to your success.




\end{document}
